% GENERATED FILE. Edit canonical lessons, then run npm run generate:books.
% canonical-source-hash: fnv1a64:eaff482722cc7c86
% canonical-lessons: MR-C531-phalandaj, MR-C531-golandaj, MR-C531-dhavpatu, MR-C531-tolavne, MR-C531-jhelne

\chapter{Players and Plays}
\label{ch:mr-players-and-plays}
\hlchaptermodality{531}

\begin{chapteropening}
\textbf{By the end of this chapter:} \emph{I can talk about batsmen, bowlers, runners and catches.}

Complete the last lesson of chapter 531: I can talk about batsmen, bowlers, runners and catches.
\end{chapteropening}

\section[\texorpdfstring{\mr{फलंदाज} (phalandāj)}{फलंदाज (phalandāj)}]{\mr{फलंदाज} (phalandāj) — a batsman; a batter}

\label{lesson:MR-C531-phalandaj}



\begin{quote}
Before the new one: say the Marathi for a contestant, then the Marathi for to carry off.
\end{quote}

\subsection*{You'll want to know: \mr{फलंदाज}}
\textbf{\mr{फलंदाज}} — \emph{phalandāj} — \textquotedblleft{}a batsman; a batter\textquotedblright{}.

\begin{quote}
\mr{फलंदाजाने} \mr{चेंडू} \mr{मैदानाबाहेर} \mr{मारला}.
\end{quote}

\begin{quote}
\emph{phalandājāne ceṇḍū maidānābāher mārlā.}
\end{quote}

\begin{quote}
The batsman hit the ball out of the ground.
\end{quote}

Masculine.

\begin{cousinweb}
\textbf{\mr{फलंदाज}} (\emph{phalandāj}) is a Marathi coinage made on the pattern of \mr{गोलंदाज}, the bowler, whose -\mr{अंदाज} is Persian \emph{andāz}, \textquotedblleft{}thrower\textquotedblright{}. The first part is probably Sanskrit \emph{phala}, glossed also as a blade (Monier-Williams): the flat face of the bat. That reading is not settled.
\end{cousinweb}

\begin{grammarlens}[title={What you've built}]
One more word for describing what you lived and felt.
\end{grammarlens}

\subsection*{Your turn}
\begin{itemize}
\raggedright
  \item \emph{Say it:} \emph{phalandāj}
  \item \emph{Say it:} \emph{phalandāj}, once more
  \item \emph{Say it:} \emph{phalandājāne ceṇḍū maidānābāher mārlā.} — The batsman hit the ball out of the ground.
  \item \emph{Recall:} say the Marathi for a contestant, then the Marathi for to carry off, then say \emph{phalandāj} again
\end{itemize}

\begin{tcolorbox}[breakable,colback=teal!4,colframe=teal!35!black,title={Before you move on}]
What does \mr{फलंदाज} mean? (\textquotedblleft{}A batsman; a batter\textquotedblright{}.) Say it once more.
\end{tcolorbox}
\section[\texorpdfstring{\mr{गोलंदाज} (golandāj)}{गोलंदाज (golandāj)}]{\mr{गोलंदाज} (golandāj) — a bowler}

\label{lesson:MR-C531-golandaj}



\begin{quote}
Before the new one: say the Marathi for to carry off, then the Marathi for a batsman.
\end{quote}

\subsection*{You'll want to know: \mr{गोलंदाज}}
\textbf{\mr{गोलंदाज}} — \emph{golandāj} — \textquotedblleft{}a bowler\textquotedblright{}.

\begin{quote}
\mr{गोलंदाजाने} \mr{खूप} \mr{वेगाने} \mr{चेंडू} \mr{टाकला}.
\end{quote}

\begin{quote}
\emph{golandājāne khūp vegāne ceṇḍū ṭāklā.}
\end{quote}

\begin{quote}
The bowler bowled the ball very fast.
\end{quote}

Masculine.

\begin{cousinweb}
\textbf{\mr{गोलंदाज}} (\emph{golandāj}) is a loan from Persian \emph{golandāz}, \textquotedblleft{}gunner\textquotedblright{}, from \emph{gola}, \textquotedblleft{}a ball\textquotedblright{}, and \emph{andāz}, \textquotedblleft{}thrower\textquotedblright{}.
\end{cousinweb}

\begin{grammarlens}[title={What you've built}]
One more word for describing what you lived and felt.
\end{grammarlens}

\subsection*{Your turn}
\begin{itemize}
\raggedright
  \item \emph{Say it:} \emph{golandāj}
  \item \emph{Say it:} \emph{golandāj}, once more
  \item \emph{Say it:} \emph{golandājāne khūp vegāne ceṇḍū ṭāklā.} — The bowler bowled the ball very fast.
  \item \emph{Recall:} say the Marathi for to carry off, then the Marathi for a batsman, then say \emph{golandāj} again
\end{itemize}

\begin{tcolorbox}[breakable,colback=teal!4,colframe=teal!35!black,title={Before you move on}]
What does \mr{गोलंदाज} mean? (\textquotedblleft{}A bowler\textquotedblright{}.) Say it once more.
\end{tcolorbox}
\section[\texorpdfstring{\mr{धावपटू} (dhāvpaṭū)}{धावपटू (dhāvpaṭū)}]{\mr{धावपटू} (dhāvpaṭū) — a runner; an athlete}

\label{lesson:MR-C531-dhavpatu}



\begin{quote}
Before the new one: say the Marathi for a batsman, then the Marathi for a bowler.
\end{quote}

\subsection*{You'll want to know: \mr{धावपटू}}
\textbf{\mr{धावपटू}} — \emph{dhāvpaṭū} — \textquotedblleft{}a runner; an athlete\textquotedblright{}.

\begin{quote}
\mr{तो} \mr{धावपटू} \mr{रोज} \mr{सकाळी} \mr{मैदानावर} \mr{धावतो}.
\end{quote}

\begin{quote}
\emph{to dhāvpaṭū roj sakāḷī maidānāvar dhāvto.}
\end{quote}

\begin{quote}
That runner runs on the ground every morning.
\end{quote}

Masculine.

\begin{cousinweb}
\textbf{\mr{धावपटू}} (\emph{dhāvpaṭū}) is built from \mr{धाव}, \textquotedblleft{}a run\textquotedblright{}, from \mr{धावणे}, and the Sanskrit \emph{paṭu}, \textquotedblleft{}skilful\textquotedblright{} (Monier-Williams).
\end{cousinweb}

\begin{grammarlens}[title={What you've built}]
One more word for describing what you lived and felt.
\end{grammarlens}

\subsection*{Your turn}
\begin{itemize}
\raggedright
  \item \emph{Say it:} \emph{dhāvpaṭū}
  \item \emph{Say it:} \emph{dhāvpaṭū}, once more
  \item \emph{Say it:} \emph{to dhāvpaṭū roj sakāḷī maidānāvar dhāvto.} — That runner runs on the ground every morning.
  \item \emph{Recall:} say the Marathi for a batsman, then the Marathi for a bowler, then say \emph{dhāvpaṭū} again
\end{itemize}

\begin{tcolorbox}[breakable,colback=teal!4,colframe=teal!35!black,title={Before you move on}]
What does \mr{धावपटू} mean? (\textquotedblleft{}A runner; an athlete\textquotedblright{}.) Say it once more.
\end{tcolorbox}
\section[\texorpdfstring{\mr{टोलवणे} (ṭolavṇe)}{टोलवणे (ṭolavṇe)}]{\mr{टोलवणे} (ṭolavṇe) — to hit away; to fend off}

\label{lesson:MR-C531-tolavne}



\begin{quote}
Before the new one: say the Marathi for a bowler, then the Marathi for a runner.
\end{quote}

\subsection*{You'll want to know: \mr{टोलवणे}}
\textbf{\mr{टोलवणे}} — \emph{ṭolavṇe} — \textquotedblleft{}to hit away; to fend off\textquotedblright{}.

\begin{quote}
\mr{फलंदाजाने} \mr{पहिला} \mr{चेंडू} \mr{लांब} \mr{टोलवला}.
\end{quote}

\begin{quote}
\emph{phalandājāne pahilā ceṇḍū lāmb ṭolavlā.}
\end{quote}

\begin{quote}
The batsman hit the first ball a long way.
\end{quote}

Past, with -\mr{ने} (he …): \mr{त्याने} \mr{चेंडू} \mr{टोलवला}, \emph{tyāne ceṇḍū ṭolavlā}.

\begin{cousinweb}
\textbf{\mr{टोलवणे}} (\emph{ṭolavṇe}) is a Marathi verb for knocking something away; its origin is not settled.
\end{cousinweb}

\begin{grammarlens}[title={What you've built}]
One more word for describing what you lived and felt.
\end{grammarlens}

\subsection*{Your turn}
\begin{itemize}
\raggedright
  \item \emph{Say it:} \emph{ṭolavṇe}
  \item \emph{Say it:} \emph{ṭolavṇe}, once more
  \item \emph{Say it:} \emph{phalandājāne pahilā ceṇḍū lāmb ṭolavlā.} — The batsman hit the first ball a long way.
  \item \emph{Recall:} say the Marathi for a bowler, then the Marathi for a runner, then say \emph{ṭolavṇe} again
\end{itemize}

\begin{tcolorbox}[breakable,colback=teal!4,colframe=teal!35!black,title={Before you move on}]
What does \mr{टोलवणे} mean? (\textquotedblleft{}To hit away; to fend off\textquotedblright{}.) Say it once more.
\end{tcolorbox}
\section[\texorpdfstring{\mr{झेलणे} (jhelṇe)}{झेलणे (jhelṇe)}]{\mr{झेलणे} (jhelṇe) — to catch a ball}

\label{lesson:MR-C531-jhelne}



\begin{quote}
Before the new one: say the Marathi for a runner, then the Marathi for to hit away.
\end{quote}

\subsection*{You'll want to know: \mr{झेलणे}}
\textbf{\mr{झेलणे}} — \emph{jhelṇe} — \textquotedblleft{}to catch a ball\textquotedblright{}.

\begin{quote}
\mr{मीराने} \mr{एका} \mr{हाताने} \mr{चेंडू} \mr{झेलला}.
\end{quote}

\begin{quote}
\emph{mīrāne ekā hātāne ceṇḍū jhellā.}
\end{quote}

\begin{quote}
Meera caught the ball with one hand.
\end{quote}

Past, with -\mr{ने} (he …): \mr{त्याने} \mr{चेंडू} \mr{झेलला}, \emph{tyāne ceṇḍū jhellā}.

\begin{cousinweb}
\textbf{\mr{झेलणे}} (\emph{jhelṇe}) is inherited: CDIAL lists Marathi \emph{jhelṇẽ}, \textquotedblleft{}to catch\textquotedblright{}, under \emph{*jhēlati}, \textquotedblleft{}undergoes\textquotedblright{} (CDIAL 5413).
\end{cousinweb}

\begin{grammarlens}[title={What you've built}]
One more word for describing what you lived and felt.
\end{grammarlens}

\subsection*{Your turn}
\begin{itemize}
\raggedright
  \item \emph{Say it:} \emph{jhelṇe}
  \item \emph{Say it:} \emph{jhelṇe}, once more
  \item \emph{Say it:} \emph{mīrāne ekā hātāne ceṇḍū jhellā.} — Meera caught the ball with one hand.
  \item \emph{Recall:} say the Marathi for a runner, then the Marathi for to hit away, then say \emph{jhelṇe} again
\end{itemize}

\begin{tcolorbox}[breakable,colback=teal!4,colframe=teal!35!black,title={Before you move on}]
What does \mr{झेलणे} mean? (\textquotedblleft{}To catch a ball\textquotedblright{}.) Say it once more.
\end{tcolorbox}
